% GENERATED FILE. Edit canonical lessons, then run npm run generate:books.
% canonical-source-hash: fnv1a64:c277ec60087059ec
% canonical-lessons: IT-C25-naso, IT-C25-stomaco, IT-C25-gola

\chapter{Nose, Stomach, and Throat}
\label{ch:it-nose-stomach-throat}
\hlchaptermodality{25}

\begin{chapteropening}
\textbf{By the end of this chapter:} \emph{I can name my nose, stomach and throat in Italian, closing the body-word set Chapter 17 opened.}

The chapter's last lesson gathers il naso, lo stomaco and la gola, closing the seven-word body set begun in Chapter 24, and lays out the five different kinds of word-history behind those seven words in one place.
\end{chapteropening}

\section[il naso]{il naso — the nose, and a flower named after it}

\label{lesson:IT-C25-naso}



\begin{quote}
Three body words left. This one has a straightforward Latin root and a genuinely surprising English cousin waiting at the end of it.
\end{quote}

\subsection*{You'll want to know: the word}
\begin{quote}
\textbf{il naso} — the nose. Masculine.
\end{quote}

Said \emph{NAH-zoh}: the single \textbf{s} sits between two vowels, which Italian voices — the same sound as the \textbf{s} in English \emph{rose}, not the one in \emph{sink}.

\begin{cousinweb}
\textbf{naso} $\leftarrow$ Latin \textbf{nasus}. English kept it directly in \textbf{nasal}, and built one word by welding it to Old English: \textbf{nostril}, \textquotedblleft{}nose\textquotedblright{} plus \emph{thyrl}, \textquotedblleft{}hole\textquotedblright{} — a nose-hole. The real surprise is a garden flower: \textbf{nasturtium} $\leftarrow$ \emph{nasus} plus \emph{torquēre}, \textquotedblleft{}to twist\textquotedblright{} — literally a \textbf{nose-twister}, for the sharp, peppery smell that makes a nose wrinkle.
\end{cousinweb}

\begin{grammarlens}[title={Grammar Lens: when a single \emph{s} sounds like \emph{z}}]
Italian's traps list already warned you that \textbf{z} can be \emph{ts} or \emph{dz}. \textbf{S} has its own split: a single \textbf{s} between two vowels is usually \textbf{voiced}, like English \emph{z} — \emph{naso}, \emph{così} (Chapter 2). A \textbf{doubled} \emph{ss} stays voiceless, the plain \emph{s} of English \emph{sink}. One letter, two sounds, and position is what decides.
\end{grammarlens}

\subsection*{Your turn}
Say these aloud:

\begin{itemize}
\raggedright
  \item \textquotedblleft{}il naso\textquotedblright{} — \emph{NAH-zoh}, voiced \emph{s}
  \item \textquotedblleft{}naso\textquotedblright{} then English \textquotedblleft{}nasal, nostril, nasturtium\textquotedblright{}
  \item the voiced-\emph{s} pair — \textquotedblleft{}naso, così\textquotedblright{}
\end{itemize}

\begin{tcolorbox}[breakable,colback=teal!4,colframe=teal!35!black,title={Before you move on}]
Give \textquotedblleft{}the nose.\textquotedblright{} (\textbf{Il naso}.) What garden flower hides \emph{nasus} plus a Latin verb for \textquotedblleft{}to twist,\textquotedblright{} and what does it literally mean? (\textbf{Nasturtium} — \textquotedblleft{}nose-twister.\textquotedblright{}) What English word welds \emph{nasus} to an Old English word for \textquotedblleft{}hole\textquotedblright{}? (\textbf{Nostril}.) How does a single \emph{s} between two vowels sound in Italian? (\textbf{Voiced}, like English \emph{z} — as in \emph{naso} or \emph{così}.)
\end{tcolorbox}
\section[lo stomaco]{lo stomaco — the stomach, on loan from Greek}

\label{lesson:IT-C25-stomaco}



\begin{quote}
Every body word so far went straight back to Latin, and the last one taught you that a single \emph{s} between vowels is voiced. This one takes a detour through Greek instead — and needs Chapter 22's other article again.
\end{quote}

\subsection*{You'll want to know: the word}
\begin{quote}
\textbf{lo stomaco} — the stomach. Masculine, and \textbf{lo}, not \emph{il} — it starts with \textbf{s} plus a consonant, exactly the cluster Chapter 22's \emph{lo zucchero} trained you to expect \emph{lo} before.
\end{quote}

\begin{cousinweb}
\textbf{stomaco} $\leftarrow$ Latin \textbf{stomachus} $\leftarrow$ Greek \textbf{stómakhos}, \textquotedblleft{}throat, gullet, stomach,\textquotedblright{} itself built on \textbf{stóma}, \textquotedblleft{}mouth.\textquotedblright{} Latin borrowed the Greek medical vocabulary wholesale in antiquity, long before Italian existed, so \emph{stomaco} only \emph{feels} as native as \emph{cuore} or \emph{naso} — underneath, it travelled through one more language than they did. English took the same word by the same route: \textbf{stomach}, spelled almost as Latin left it.
\end{cousinweb}

\subsection*{Your turn}
Say these aloud:

\begin{itemize}
\raggedright
  \item \textquotedblleft{}lo stomaco\textquotedblright{} — \emph{lo}, not \emph{il}
  \item \textquotedblleft{}stomaco\textquotedblright{} then English \textquotedblleft{}stomach\textquotedblright{} — nearly the same word
  \item the \emph{lo} pair — \textquotedblleft{}lo zucchero, lo stomaco\textquotedblright{}
\end{itemize}

\begin{tcolorbox}[breakable,colback=teal!4,colframe=teal!35!black,title={Before you move on}]
Give \textquotedblleft{}the stomach.\textquotedblright{} (\textbf{Lo stomaco}.) Why \emph{lo} and not \emph{il}? (It starts with \textbf{s} plus a consonant.) Which language did \emph{stomaco} pass through before it reached Latin? (\textbf{Greek} — \emph{stómakhos}, from \emph{stóma}, \textquotedblleft{}mouth.\textquotedblright{}) Is \emph{stomaco}'s root older or younger, inside Latin, than \emph{naso}'s? (\textbf{Younger} — borrowed from Greek, where \emph{nasus} was native Latin all along.)
\end{tcolorbox}
\section[la gola]{la gola — the throat, closing the body}

\label{lesson:IT-C25-gola}



\begin{quote}
The last body word, and the piece that connects the mouth to the stomach — heart, eyes, ears, mouth, nose, stomach, and now the throat between them.
\end{quote}

\subsection*{You'll want to know: the word}
\begin{quote}
\textbf{la gola} — the throat. Feminine, regular \textbf{-a}.
\end{quote}

\begin{cousinweb}
\textbf{gola} $\leftarrow$ Latin \textbf{gula}, \textquotedblleft{}throat.\textquotedblright{} English kept it almost whole in \textbf{gullet}, by way of Old French. A tempting further step: Latin also had a verb \textbf{gluttire}, \textquotedblleft{}to gulp down,\textquotedblright{} the direct ancestor of English \textbf{glutton}. Several dictionaries propose that \emph{gluttire} and \emph{gula} ultimately share the same ancient root for swallowing — plausible, and worth knowing, but treated here as \textbf{proposed}, not settled, exactly as this book flagged \emph{persona}'s Etruscan origin as the leading guess rather than a certainty.
\end{cousinweb}

\begin{grammarlens}[title={What you've built this chapter}]
Two chapters, seven words: \textbf{cuore, occhio, orecchio, bocca, naso, stomaco, gola} — heart, eyes, ears, mouth, nose, stomach, throat. Three came straight from Latin roots (\emph{cuore}, \emph{naso}, \emph{gola}); two collapsed under the same Italian sound-law (\emph{occhio}, \emph{orecchio}); one is slang that outlived the \textquotedblleft{}proper\textquotedblright{} word it replaced (\emph{bocca}); one detoured through Greek before it ever reached Latin (\emph{stomaco}). Seven words, five different kinds of history, all landing on one \emph{Come stai?}.
\end{grammarlens}

\subsection*{Your turn}
Say these aloud:

\begin{itemize}
\raggedright
  \item \textquotedblleft{}la gola\textquotedblright{}
  \item \textquotedblleft{}gola\textquotedblright{} then English \textquotedblleft{}gullet\textquotedblright{} — and, cautiously, \textquotedblleft{}glutton\textquotedblright{}
  \item the full seven — \textquotedblleft{}cuore, occhio, orecchio, bocca, naso, stomaco, gola\textquotedblright{}
\end{itemize}

\begin{tcolorbox}[breakable,colback=teal!4,colframe=teal!35!black,title={Before you move on}]
Give \textquotedblleft{}the throat.\textquotedblright{} (\textbf{La gola}.) What Latin word gave Italian \emph{gola}, and its clear English cousin? (\emph{Gula} $\to$ \textbf{gullet}.) Which English word is \emph{proposed}, not certain, to share the same root? (\textbf{Glutton}, from \emph{gluttire}.) Name the seven body words this pair of chapters taught, in order. (\emph{Cuore, occhio, orecchio, bocca, naso, stomaco, gola}.)
\end{tcolorbox}
